
\documentclass[12pt]{article}

\usepackage[margin=1in]{geometry}
\usepackage{amsmath,amssymb}
\usepackage{booktabs}
\usepackage{array}
\usepackage{longtable}
\usepackage{hyperref}
\usepackage{siunitx}
\usepackage{enumitem}
\usepackage{setspace}

\onehalfspacing

\title{Waves, Coherence, Lasers, Modulation, Quantum Wave Phenomena, and Radiation Shielding}
\author{Zachary Tullier}
\date{6/24/26}

\begin{document}

\maketitle
\tableofcontents
\newpage

\section{Identification Parameters of Waves and Their Quantum-Mechanical Relationships}

A wave is a disturbance that transfers energy and momentum through space or matter. A monochromatic harmonic wave traveling in the positive $x$-direction may be written as

\begin{equation}
\psi(x,t)=A\cos(kx-\omega t+\phi_0),
\end{equation}

where $A$ is the amplitude, $k$ is the wave number, $\omega$ is the angular frequency, and $\phi_0$ is the initial phase.

\subsection{Wavelength}

The wavelength $\lambda$ is the spatial distance over which the wave repeats itself:

\begin{equation}
\psi(x+\lambda,t)=\psi(x,t).
\end{equation}

The wave number is related to wavelength by

\begin{equation}
k=\frac{2\pi}{\lambda}.
\end{equation}

For electromagnetic radiation in vacuum,

\begin{equation}
c=\lambda \nu,
\end{equation}

where $c$ is the speed of light and $\nu$ is the frequency. Therefore,

\begin{equation}
\lambda=\frac{c}{\nu}.
\end{equation}

Quantum mechanically, a photon has momentum

\begin{equation}
p=\hbar k=\frac{h}{\lambda},
\end{equation}

which is the de Broglie relation. Thus,

\begin{equation}
\boxed{\lambda=\frac{h}{p}}.
\end{equation}

A shorter wavelength corresponds to a larger photon momentum.

\subsection{Frequency}

Frequency $\nu$ is the number of complete oscillations per unit time:

\begin{equation}
\nu=\frac{N}{t}.
\end{equation}

Its SI unit is the hertz:

\begin{equation}
1~\mathrm{Hz}=1~\mathrm{s}^{-1}.
\end{equation}

Angular frequency is

\begin{equation}
\omega=2\pi\nu.
\end{equation}

Quantum mechanically, the energy of one photon is

\begin{equation}
\boxed{E=h\nu=\hbar\omega}.
\end{equation}

Combining $c=\lambda\nu$ with $E=h\nu$ gives

\begin{equation}
\boxed{E=\frac{hc}{\lambda}}.
\end{equation}

Thus, photon energy is inversely proportional to wavelength.

\subsection{Period}

The period $T$ is the time required for one complete oscillation:

\begin{equation}
T=\frac{1}{\nu}.
\end{equation}

Since $E=h\nu$,

\begin{equation}
E=\frac{h}{T},
\end{equation}

and therefore,

\begin{equation}
\boxed{T=\frac{h}{E}}.
\end{equation}

The associated quantum state evolves with angular frequency

\begin{equation}
\omega=\frac{E}{\hbar}.
\end{equation}

The time dependence of an energy eigenstate is

\begin{equation}
\Psi(\mathbf{r},t)=\psi(\mathbf{r})e^{-iEt/\hbar}.
\end{equation}

\subsection{Energy}

For electromagnetic radiation, the energy density is

\begin{equation}
u=\frac{1}{2}\epsilon_0 E_0^2+\frac{1}{2\mu_0}B_0^2.
\end{equation}

For a plane electromagnetic wave,

\begin{equation}
B_0=\frac{E_0}{c}.
\end{equation}

The electric and magnetic contributions are equal, so the average total energy density is

\begin{equation}
\langle u\rangle=\frac{1}{2}\epsilon_0 E_0^2.
\end{equation}

The time-averaged intensity is

\begin{equation}
\boxed{I=\frac{1}{2}c\epsilon_0 E_0^2}.
\end{equation}

Therefore,

\begin{equation}
I\propto E_0^2.
\end{equation}

Quantum mechanically, if $N$ photons of frequency $\nu$ pass through an area $A_s$ during time $\Delta t$, then

\begin{equation}
I=\frac{Nh\nu}{A_s\Delta t}.
\end{equation}

\subsection{Amplitude in Terms of Energy}

For a classical mechanical wave,

\begin{equation}
E_{\text{wave}}\propto A^2.
\end{equation}

For a harmonic oscillator,

\begin{equation}
E=\frac{1}{2}m\omega^2A^2.
\end{equation}

Therefore,

\begin{equation}
A=\sqrt{\frac{2E}{m\omega^2}}.
\end{equation}

For an electromagnetic wave occupying a volume $V$,

\begin{equation}
U=\langle u\rangle V=\frac{1}{2}\epsilon_0E_0^2V.
\end{equation}

Hence,

\begin{equation}
\boxed{E_0=\sqrt{\frac{2U}{\epsilon_0V}}}.
\end{equation}

If the electromagnetic field contains $N$ photons,

\begin{equation}
U=Nh\nu,
\end{equation}

so that

\begin{equation}
\boxed{E_0=\sqrt{\frac{2Nh\nu}{\epsilon_0V}}}.
\end{equation}

Thus,

\begin{equation}
E_0\propto \sqrt{N}.
\end{equation}

For a particle wavefunction $\Psi$,

\begin{equation}
\boxed{|\Psi(\mathbf r,t)|^2}
\end{equation}

is the probability density of finding the particle at position $\mathbf r$ and time $t$.

\subsection{Summary of Quantum Relations}

The principal wave--particle relations are

\begin{equation}
\boxed{E=h\nu=\hbar\omega},
\end{equation}

\begin{equation}
\boxed{p=\frac{h}{\lambda}=\hbar k},
\end{equation}

\begin{equation}
\boxed{c=\lambda\nu},
\end{equation}

and for a photon,

\begin{equation}
\boxed{E=pc}.
\end{equation}

Therefore,

\begin{equation}
E=\frac{hc}{\lambda},
\qquad
p=\frac{E}{c}.
\end{equation}

\section{Phase and Coherence}

\subsection{Phase of a Wave}

The phase of the harmonic wave

\begin{equation}
\psi(x,t)=A\cos(kx-\omega t+\phi_0)
\end{equation}

is

\begin{equation}
\boxed{\phi(x,t)=kx-\omega t+\phi_0}.
\end{equation}

For two waves,

\begin{align}
\psi_1 &= A_1\cos(kx-\omega t+\phi_1),\\
\psi_2 &= A_2\cos(kx-\omega t+\phi_2),
\end{align}

the phase difference is

\begin{equation}
\Delta\phi=\phi_2-\phi_1.
\end{equation}

Constructive interference occurs when

\begin{equation}
\Delta\phi=2m\pi,
\end{equation}

whereas destructive interference occurs when

\begin{equation}
\Delta\phi=(2m+1)\pi.
\end{equation}

The resultant intensity is

\begin{equation}
\boxed{
I=I_1+I_2+2\sqrt{I_1I_2}\cos\Delta\phi
}.
\end{equation}

For equal intensities $I_1=I_2=I_0$,

\begin{equation}
I=2I_0(1+\cos\Delta\phi)
=4I_0\cos^2\left(\frac{\Delta\phi}{2}\right).
\end{equation}

\subsection{Coherence}

Coherence describes the degree to which a wave maintains a predictable phase relationship. Two waves are perfectly coherent if their phase difference remains constant:

\begin{equation}
\frac{d}{dt}\Delta\phi=0.
\end{equation}

The normalized first-order coherence function is

\begin{equation}
g^{(1)}(\mathbf r_1,t_1;\mathbf r_2,t_2)
=
\frac{
\langle E^*(\mathbf r_1,t_1)E(\mathbf r_2,t_2)\rangle
}{
\sqrt{
\langle |E(\mathbf r_1,t_1)|^2\rangle
\langle |E(\mathbf r_2,t_2)|^2\rangle
}
}.
\end{equation}

Perfect coherence corresponds to

\begin{equation}
|g^{(1)}|=1,
\end{equation}

whereas complete incoherence corresponds approximately to

\begin{equation}
|g^{(1)}|=0.
\end{equation}

\subsection{Temporal Coherence}

Temporal coherence describes the correlation of a wave with itself at different times at the same spatial location. It is related to spectral bandwidth $\Delta\nu$ by

\begin{equation}
\boxed{\tau_c\sim\frac{1}{\Delta\nu}}.
\end{equation}

The coherence length is

\begin{equation}
\boxed{L_c=c\tau_c\sim\frac{c}{\Delta\nu}}.
\end{equation}

For a narrow wavelength bandwidth,

\begin{equation}
\Delta\nu\approx\frac{c\Delta\lambda}{\lambda^2}.
\end{equation}

Therefore,

\begin{equation}
\boxed{L_c\approx\frac{\lambda^2}{\Delta\lambda}}.
\end{equation}

\subsection{Spatial or Positional Coherence}

Spatial coherence describes the phase correlation between different points across the same wavefront at the same time. For an incoherent source of diameter $d$ observed at distance $R$, the transverse coherence width is approximately

\begin{equation}
\ell_{\perp}\sim\frac{\lambda R}{d}.
\end{equation}

A smaller source therefore produces greater spatial coherence.

\subsection{Best Ways to Acquire Coherence}

Highly coherent light is most effectively produced with a laser because stimulated emission generates photons with closely matched frequency, phase, direction, and polarization.

Temporal coherence can be improved using optical resonators, narrow-band interference filters, monochromators, and single-longitudinal-mode operation. Spatial coherence can be improved with pinholes, single-mode optical fibers, spatial filters, and well-aligned optical cavities.

\section{LASER: Meaning, Production, and Distinctive Features}

\subsection{Meaning of LASER}

LASER stands for

\begin{equation}
\boxed{\text{Light Amplification by Stimulated Emission of Radiation}}.
\end{equation}

\subsection{Atomic Processes Involved}

Consider two atomic energy levels $E_1$ and $E_2$, where

\begin{equation}
E_2-E_1=h\nu.
\end{equation}

\subsubsection{Absorption}

\begin{equation}
\mathrm{Atom}(E_1)+h\nu
\rightarrow
\mathrm{Atom}(E_2).
\end{equation}

\subsubsection{Spontaneous Emission}

\begin{equation}
\mathrm{Atom}(E_2)
\rightarrow
\mathrm{Atom}(E_1)+h\nu.
\end{equation}

\subsubsection{Stimulated Emission}

\begin{equation}
\mathrm{Atom}(E_2)+h\nu
\rightarrow
\mathrm{Atom}(E_1)+2h\nu.
\end{equation}

The stimulated photon has the same frequency, phase, direction, and polarization as the incident photon.

\subsection{Population Inversion}

At thermal equilibrium,

\begin{equation}
N_1>N_2.
\end{equation}

Laser operation requires population inversion:

\begin{equation}
\boxed{N_2>N_1}.
\end{equation}

The thermal-equilibrium population ratio is

\begin{equation}
\frac{N_2}{N_1}
=
\exp\left[-\frac{E_2-E_1}{k_BT}\right].
\end{equation}

\subsection{Optical Resonator}

A laser normally contains a gain medium, a pump source, and two mirrors forming an optical cavity.

The threshold gain condition is approximately

\begin{equation}
\boxed{
g_{\mathrm{th}}
=
\alpha_i
+
\frac{1}{2L}
\ln\left(\frac{1}{R_1R_2}\right)
},
\end{equation}

where $\alpha_i$ is the internal loss coefficient, $L$ is cavity length, and $R_1$ and $R_2$ are mirror reflectivities.

The resonant longitudinal modes satisfy

\begin{equation}
2L=m\lambda.
\end{equation}

Equivalently,

\begin{equation}
\nu_m=\frac{mc}{2nL}.
\end{equation}

\subsection{Distinctive Features of Laser Light}

Laser light is characterized by high monochromaticity, coherence, directionality, intensity, and often strong polarization.

For a diffraction-limited circular aperture,

\begin{equation}
\theta\approx1.22\frac{\lambda}{D}.
\end{equation}

For a Gaussian beam,

\begin{equation}
\theta\approx\frac{\lambda}{\pi w_0}.
\end{equation}

The intensity is

\begin{equation}
I=\frac{P}{A}.
\end{equation}

The most distinctive feature is coherent amplification by stimulated emission.

\section{Collective Behavior of Radiation Beams}

\subsection{Reflection}

Reflection obeys

\begin{equation}
\boxed{\theta_i=\theta_r}.
\end{equation}

At normal incidence, the electric-field reflection coefficient is

\begin{equation}
r=\frac{n_1-n_2}{n_1+n_2}.
\end{equation}

The reflected intensity fraction is

\begin{equation}
\boxed{
R=|r|^2
=
\left(\frac{n_1-n_2}{n_1+n_2}\right)^2
}.
\end{equation}

Quantum mechanically, an incident photon state may become

\begin{equation}
|\psi\rangle
=
r|\text{reflected}\rangle
+
t|\text{transmitted}\rangle.
\end{equation}

The associated probabilities are

\begin{equation}
P_R=|r|^2,
\qquad
P_T=|t|^2.
\end{equation}

\subsection{Refraction}

The refractive index is

\begin{equation}
n=\frac{c}{v}.
\end{equation}

Snell's law is

\begin{equation}
\boxed{n_1\sin\theta_1=n_2\sin\theta_2}.
\end{equation}

Quantum mechanically, conservation of momentum parallel to the interface gives

\begin{equation}
p_{1,\parallel}=p_{2,\parallel},
\end{equation}

or

\begin{equation}
\hbar k_1\sin\theta_1
=
\hbar k_2\sin\theta_2.
\end{equation}

Since

\begin{equation}
k=\frac{n\omega}{c},
\end{equation}

Snell's law follows.

\subsection{Transmission}

For a nonabsorbing interface,

\begin{equation}
R+T=1.
\end{equation}

With absorption,

\begin{equation}
\boxed{R+T+A=1}.
\end{equation}

For a homogeneous absorbing material,

\begin{equation}
\boxed{I(x)=I_0e^{-\mu x}}.
\end{equation}

Quantum mechanically, the survival probability is

\begin{equation}
P_{\text{survival}}(x)=e^{-\mu x},
\end{equation}

and the interaction probability is

\begin{equation}
P_{\text{interaction}}(x)=1-e^{-\mu x}.
\end{equation}

\section{Modulation: Classical or Quantum?}

Modulation is primarily a classical wave phenomenon, although it also has a quantum description. A carrier wave may be written as

\begin{equation}
c(t)=A_c\cos(2\pi f_ct+\phi_c).
\end{equation}

\subsection{Amplitude Modulation}

For a message

\begin{equation}
m(t)=A_m\cos(2\pi f_mt),
\end{equation}

the amplitude-modulated signal is

\begin{equation}
\boxed{
s_{\text{AM}}(t)
=
A_c
\left[
1+m_a\cos(2\pi f_mt)
\right]
\cos(2\pi f_ct)
},
\end{equation}

where

\begin{equation}
m_a=\frac{A_m}{A_c}.
\end{equation}

Expanding,

\begin{align}
s_{\text{AM}}(t)
&=
A_c\cos(2\pi f_ct)
+
\frac{m_aA_c}{2}
\cos[2\pi(f_c+f_m)t] \nonumber\\
&\quad+
\frac{m_aA_c}{2}
\cos[2\pi(f_c-f_m)t].
\end{align}

\subsection{Frequency Modulation}

The instantaneous frequency is

\begin{equation}
f_i(t)=f_c+k_fm(t).
\end{equation}

The phase is

\begin{equation}
\phi(t)
=
2\pi\int_0^t f_i(t')\,dt'.
\end{equation}

For sinusoidal modulation,

\begin{equation}
\boxed{
s_{\text{FM}}(t)
=
A_c
\cos
\left[
2\pi f_ct+\beta\sin(2\pi f_mt)
\right]
},
\end{equation}

where

\begin{equation}
\beta=\frac{\Delta f}{f_m}.
\end{equation}

The FM spectrum can be written as

\begin{equation}
s_{\text{FM}}(t)
=
A_c
\sum_{n=-\infty}^{\infty}
J_n(\beta)
\cos[2\pi(f_c+nf_m)t].
\end{equation}

\subsection{Phase Modulation}

\begin{equation}
\boxed{
s_{\text{PM}}(t)
=
A_c\cos[2\pi f_ct+k_pm(t)]
}.
\end{equation}

Instantaneous frequency is related to phase by

\begin{equation}
f_i(t)
=
\frac{1}{2\pi}\frac{d\phi}{dt}.
\end{equation}

\subsection{Pulse Modulation}

Pulse modulation includes pulse amplitude modulation, pulse width modulation, pulse position modulation, and pulse code modulation.

A pulse of duration $\Delta t$ has a finite bandwidth satisfying approximately

\begin{equation}
\Delta\nu\,\Delta t\gtrsim\frac{1}{4\pi}.
\end{equation}

\subsection{Quantum Modulation}

Information may be encoded in photon number, phase, polarization, frequency, arrival time, or orbital angular momentum.

For polarization encoding,

\begin{equation}
|0\rangle=|H\rangle,
\qquad
|1\rangle=|V\rangle.
\end{equation}

A general superposition is

\begin{equation}
|\psi\rangle
=
\alpha|H\rangle+\beta|V\rangle,
\end{equation}

where

\begin{equation}
|\alpha|^2+|\beta|^2=1.
\end{equation}

\section{Quantum-Mechanical Wave Phenomena Observable Classically}

\subsection{Interference}

Quantum amplitudes add:

\begin{equation}
\Psi=\Psi_1+\Psi_2.
\end{equation}

The probability density becomes

\begin{equation}
|\Psi|^2
=
|\Psi_1|^2
+
|\Psi_2|^2
+
2\operatorname{Re}(\Psi_1^*\Psi_2).
\end{equation}

The cross term produces interference.

\subsection{Diffraction}

For a single slit,

\begin{equation}
I(\theta)
=
I_0
\left(
\frac{\sin\beta}{\beta}
\right)^2,
\end{equation}

where

\begin{equation}
\beta=\frac{\pi a\sin\theta}{\lambda}.
\end{equation}

The minima occur when

\begin{equation}
a\sin\theta=m\lambda,
\qquad
m=\pm1,\pm2,\ldots
\end{equation}

For matter waves,

\begin{equation}
\lambda=\frac{h}{p}.
\end{equation}

\subsection{Superposition}

Classical waves obey

\begin{equation}
\psi=\psi_1+\psi_2.
\end{equation}

Quantum states obey

\begin{equation}
|\Psi\rangle
=
c_1|\psi_1\rangle
+
c_2|\psi_2\rangle.
\end{equation}

\subsection{Tunneling Analogues}

For a rectangular barrier with $V_0>E$,

\begin{equation}
\psi(x)\propto e^{-\kappa x},
\end{equation}

where

\begin{equation}
\kappa
=
\frac{\sqrt{2m(V_0-E)}}{\hbar}.
\end{equation}

The transmission probability approximately satisfies

\begin{equation}
\boxed{T\propto e^{-2\kappa a}}.
\end{equation}

Classical analogues include evanescent electromagnetic coupling, frustrated total internal reflection, acoustic tunneling, and waveguide cutoff transmission.

\subsection{Standing Waves and Quantization}

For a string fixed at both ends,

\begin{equation}
L=\frac{n\lambda_n}{2},
\end{equation}

so that

\begin{equation}
\lambda_n=\frac{2L}{n}.
\end{equation}

The allowed frequencies are

\begin{equation}
f_n=\frac{nv}{2L}.
\end{equation}

For a quantum particle in an infinite square well,

\begin{equation}
E_n=\frac{n^2\pi^2\hbar^2}{2mL^2}.
\end{equation}

\section{Radiation Shielding and Its Quantum-Mechanical Explanation}

Radiation shielding reduces radiation intensity or energy by forcing radiation to interact with matter through discrete quantum processes.

\subsection{Attenuation Law}

For a narrow monoenergetic photon beam,

\begin{equation}
\boxed{I(x)=I_0e^{-\mu x}}.
\end{equation}

The attenuation coefficient is

\begin{equation}
\mu=n\sigma,
\end{equation}

where $n$ is the target number density and $\sigma$ is the microscopic cross section.

The mean free path is

\begin{equation}
\boxed{\ell=\frac{1}{\mu}=\frac{1}{n\sigma}}.
\end{equation}

The half-value layer is

\begin{equation}
\boxed{x_{1/2}=\frac{\ln 2}{\mu}}.
\end{equation}

The tenth-value layer is

\begin{equation}
\boxed{x_{1/10}=\frac{\ln 10}{\mu}}.
\end{equation}

\subsection{Photon Shielding}

\subsubsection{Photoelectric Effect}

\begin{equation}
\gamma+\mathrm{Atom}
\rightarrow
\mathrm{Atom}^+
+
e^-.
\end{equation}

The photoelectron kinetic energy is

\begin{equation}
\boxed{K_e=h\nu-E_b}.
\end{equation}

The approximate dependence is

\begin{equation}
\sigma_{\mathrm{pe}}\propto\frac{Z^n}{E_\gamma^3}.
\end{equation}

\subsubsection{Compton Scattering}

\begin{equation}
\gamma+e^-
\rightarrow
\gamma'+e^-.
\end{equation}

The Compton wavelength shift is

\begin{equation}
\boxed{
\lambda'-\lambda
=
\frac{h}{m_ec}(1-\cos\theta)
}.
\end{equation}

The scattered photon energy is

\begin{equation}
\boxed{
E_\gamma'
=
\frac{E_\gamma}
{
1+
\frac{E_\gamma}{m_ec^2}(1-\cos\theta)
}
}.
\end{equation}

\subsubsection{Pair Production}

Pair production becomes possible when

\begin{equation}
E_\gamma\geq 2m_ec^2=1.022~\mathrm{MeV}.
\end{equation}

The reaction is

\begin{equation}
\gamma+Z
\rightarrow
e^-+e^++Z.
\end{equation}

The kinetic energy balance is

\begin{equation}
K_{e^-}+K_{e^+}
=
E_\gamma-1.022~\mathrm{MeV}.
\end{equation}

Annihilation produces

\begin{equation}
e^++e^-
\rightarrow
2\gamma,
\end{equation}

typically with each photon having energy

\begin{equation}
511~\mathrm{keV}.
\end{equation}

\subsection{Charged-Particle Shielding}

The stopping power is

\begin{equation}
\boxed{
S=-\frac{dE}{dx}
}.
\end{equation}

For heavy charged particles, the Bethe relation has the approximate dependence

\begin{equation}
-\frac{dE}{dx}
\propto
\frac{z^2}{v^2}
\left[
\ln\left(\frac{2m_ev^2}{I}\right)-\cdots
\right].
\end{equation}

Ionization may be represented as

\begin{equation}
\mathrm{Atom}+q
\rightarrow
\mathrm{Atom}^++q+e^-.
\end{equation}

\subsection{Beta-Particle Shielding and Bremsstrahlung}

Bremsstrahlung production is represented by

\begin{equation}
e^-+Z
\rightarrow
e^-+Z+\gamma.
\end{equation}

A practical beta shield commonly uses a low-$Z$ material to stop electrons, followed by a high-$Z$ material to attenuate secondary photons.

\subsection{Alpha-Particle Shielding}

Alpha particles are helium nuclei:

\begin{equation}
{}^4_2\mathrm{He}^{2+}.
\end{equation}

They lose energy rapidly through ionization and can be stopped by paper, clothing, several centimeters of air, or the outer layer of skin. Internal alpha contamination is much more hazardous.

\subsection{Neutron Shielding}

Neutrons interact with nuclei by elastic scattering, inelastic scattering, capture, and nuclear reactions.

For elastic scattering, the maximum fractional energy transfer is

\begin{equation}
\boxed{
\frac{\Delta E_{\max}}{E}
=
\frac{4A}{(1+A)^2}
}.
\end{equation}

For hydrogen, $A=1$, giving

\begin{equation}
\frac{\Delta E_{\max}}{E}=1.
\end{equation}

Examples of neutron capture include

\begin{equation}
n+{}^{10}\mathrm{B}
\rightarrow
{}^7\mathrm{Li}+\alpha,
\end{equation}

and

\begin{equation}
n+{}^1\mathrm{H}
\rightarrow
{}^2\mathrm{H}+\gamma.
\end{equation}

Hydrogen-rich materials slow neutrons, while boron-containing materials absorb thermal neutrons.

\subsection{Layered Shielding}

Mixed radiation fields often require layered shielding. A neutron--gamma shield may use hydrogen-rich material for moderation, boron-containing material for neutron capture, and high-density material for secondary gamma attenuation.

Similarly, beta shielding may use low-$Z$ plastic followed by a high-$Z$ layer.

\section{Conclusion}

Classical wave quantities are connected to quantum mechanics through

\begin{equation}
E=h\nu=\hbar\omega,
\qquad
p=\frac{h}{\lambda}=\hbar k.
\end{equation}

Coherence describes stable phase relationships and may be temporal or spatial. Laser light is generated through population inversion and stimulated emission, producing highly coherent, directional, monochromatic, and intense radiation.

Reflection, refraction, and transmission follow classical electromagnetic boundary conditions and can also be described quantum mechanically through probability amplitudes, momentum conservation, and coherent scattering.

Modulation is primarily classical but can be generalized to quantum states. Interference, diffraction, superposition, resonant modes, and tunneling-like effects occur in both classical and quantum systems.

Radiation shielding works through probabilistic interactions such as photoelectric absorption, Compton scattering, pair production, ionization, nuclear scattering, and neutron capture.

\end{document}
